\newpage
\begin{center}
  \textbf{\large 3. ПРАКТИЧЕСКАЯ ЧАСТЬ}
\end{center}
\refstepcounter{chapter}
\addcontentsline{toc}{chapter}{3. ПРАКТИЧЕСКАЯ ЧАСТЬ}

Практическая часть работы основана на CFD-данных для задачи всплытия пузырька.
Численные решения получены с помощью Basilisk и используются как эталон при
оценке качества нейросетевых моделей. Для каждого расчета доступны поля
объемной доли, давления и компонент скорости:
\[
    \alpha,\qquad p,\qquad u,\qquad v.
\]
При обучении PINN напрямую используется поле \(\alpha\), задающее положение
интерфейса. Поля скорости и давления используются при сравнении предсказаний
модели с CFD-решением.

Из каждой CFD-симуляции формируются три типа обучающих точек:
\[
    \mathcal{D}_\alpha,\qquad
    \mathcal{D}_f,\qquad
    \mathcal{D}_{bc}.
\]
Множество \(\mathcal{D}_\alpha\) содержит точки для обучения объемной доли,
\(\mathcal{D}_f\) используется для вычисления физических невязок уравнений, а
\(\mathcal{D}_{bc}\) задает граничные условия. При выборе точек особое внимание
уделяется окрестности интерфейса, поскольку именно там поля меняются наиболее
резко и возникают основные ошибки аппроксимации.

Для параметрической PINN обучающие данные объединяются из нескольких расчетов с
разными начальными радиусами:
\[
    R_{\mathrm{train}}\in\{0.18,0.22,0.25,0.28,0.32\}.
\]
Для проверки способности к интерполяции по параметру используются радиусы
\[
    R_{\mathrm{test}}\in\{0.20,0.30\}.
\]
Эти радиусы лежат внутри диапазона обучающих значений, но не входят в обучающую
выборку.

Качество предсказаний оценивается относительной \(L_2\)-ошибкой. Для величины
\(q\) она определяется формулой
\begin{equation}
    L_2(q)
    =
    \frac{\sqrt{\sum_{i=1}^{N}(q_i^{\mathrm{PINN}}-q_i^{\mathrm{CFD}})^2}}
         {\sqrt{\sum_{i=1}^{N}(q_i^{\mathrm{CFD}})^2}}.
    \label{eq:chapter3_l2}
\end{equation}
Эта метрика вычисляется отдельно для \(\alpha\), \(u\), \(v\) и \(p\). Для
давления перед вычислением ошибки производится выравнивание аддитивной
константы путём согласования средних значений у верхней границы области, так
как несжимаемое давление определено с точностью до константы.

\section{Моделирование с использованием простой PINN}

\subsection{Настройка эксперимента}

Простая PINN обучается для одного фиксированного начального радиуса пузырька
\(R=0.25\). В этой постановке вход сети содержит только координаты и время:
\[
    (x^*,y^*,t^*).
\]
На выходе модель предсказывает четыре поля:
\[
    (u^*,v^*,p^*,\alpha).
\]
Радиус \(R\) не передается в сеть, поэтому модель не имеет явного механизма для
изменения начальной геометрии пузырька.

Используется архитектура с восемью скрытыми слоями по 350 нейронов и функцией
активации \(\tanh\). Обучение проводится в две стадии алгоритмом Adam:
\[
\begin{array}{c|cc}
    \hline
    \text{Стадия} & 1 & 2 \\
    \hline
    \text{Эпохи} & 5000 & 5000 \\
    \text{Шаг обучения} & 10^{-4} & 5\cdot10^{-5} \\
    \hline
\end{array}
\]
Суммарное время обучения составило 2 ч 26 мин.

\subsection{Результаты на обучающем радиусе}

Сначала простая PINN проверяется на том же радиусе \(R=0.25\), на котором
обучалась. Сравнение проводится по полям \(\alpha\), \(u\), \(v\) и \(p\)
с использованием метрики~\eqref{eq:chapter3_l2}; результаты приведены в
таблице~\ref{tab:simple_train}.

\begin{table}[H]
    \centering
    \caption{Относительные \(L_2\)-ошибки простой PINN на обучающем радиусе \(R=0.25\), момент времени \(t=3.0\).}
    \label{tab:simple_train}
    \begin{tabular}{c|cccc}
        \hline
        Поле & \(u\) & \(v\) & \(p\) & \(\alpha\) \\
        \hline
        \(L_2\), \% & 5.874 & 4.229 & 0.809 & 6.782 \\
        \hline
    \end{tabular}
\end{table}

Поле \(\alpha\) восстановлено с ошибкой около 6.8\%, скорости --- в пределах
4--6\%, давление --- менее 1\%. Наибольшая ошибка локализована в окрестности
интерфейса, где поля меняются наиболее резко. На рисунке~\ref{fig:simple_train}
приведено сравнение предсказаний PINN и CFD-решения для всех четырёх полей.

\begin{figure}[H]
    \centering
    \includegraphics[width=\textwidth]{simple_pinn_R1.00_t3.000.png}
    \caption{Простая PINN, \(R=0.25\) (обучающий), \(t=3.0\). Слева --- CFD,
    в центре --- PINN, справа --- разность. Ошибка скоростей и объемной доли
    сосредоточена вблизи интерфейса.}
    \label{fig:simple_train}
\end{figure}

\subsection{Результаты на новом радиусе}

Второй эксперимент --- проверка на радиусе \(R=0.30\), который не использовался
при обучении. Поскольку радиус не входит во входной вектор, сеть не может
получить информацию об изменившемся начальном размере пузырька. Результаты
приведены в таблице~\ref{tab:simple_test}.

\begin{table}[H]
    \centering
    \caption{Относительные \(L_2\)-ошибки простой PINN на тестовом радиусе \(R=0.30\).}
    \label{tab:simple_test}
    \begin{tabular}{c|cccc}
        \hline
        Поле & \(u\) & \(v\) & \(p\) & \(\alpha\) \\
        \hline
        \(L_2\), \% & 54.43 & 29.53 & 1.231 & 54.96 \\
        \hline
    \end{tabular}
\end{table}

Ошибки по скорости \(u\) и объемной доле \(\alpha\) выросли до 54\%,
по скорости \(v\) --- до 30\%. Только давление осталось на прежнем уровне,
поскольку оно определяется в основном гидростатической компонентой, слабо
зависящей от начального радиуса. Поля скорости и положение интерфейса,
напротив, принципиально различаются при разных \(R\), и сеть без доступа к
параметру задачи не способна их восстановить. На рисунке~\ref{fig:simple_fail}
видно, что модель предсказывает пузырёк в положении, характерном для \(R=0.25\),
тогда как реальный пузырёк \(R=0.30\) имеет другую скорость подъёма и форму.

\begin{figure}[H]
    \centering
    \includegraphics[width=\textwidth]{simple_pinn_R1.20_t3.000.png}
    \caption{Простая PINN, \(R=0.30\) (тестовый), \(t=3.0\). Модель обучена на
    \(R=0.25\). Ошибки по \(\alpha\) и скоростям недопустимо велики: интерфейс
    расположен неверно.}
    \label{fig:simple_fail}
\end{figure}

\subsection{Вывод по простой PINN}

Простая PINN хорошо восстанавливает одно течение, для которого она получала
данные. Однако перенос на другой радиус приводит к недопустимо большим ошибкам.
Чтобы одна модель могла учитывать изменение начального размера пузырька,
радиус необходимо включить во вход сети. Это приводит к параметрической PINN.

\section{Моделирование с использованием параметрической PINN}

Параметрическая PINN принимает на вход координаты, время и безразмерный
радиус \((x^*,y^*,t^*,R^*)\) и обучается одновременно на пяти CFD-симуляциях
с \(R_{\mathrm{train}}\in\{0.18,0.22,0.25,0.28,0.32\}\). Качество оценивается
на тестовых радиусах \(R\in\{0.20,0.30\}\), не входящих в обучающую выборку.
В данном разделе последовательно рассматриваются три конфигурации архитектуры,
влияние числа стадий обучения и исследование влияния числа обучающих точек.

\subsection{Лёгкая архитектура}

Лёгкая конфигурация --- наименее ёмкая из рассматриваемых: четыре скрытых
слоя по 128 нейронов с функцией активации \(\tanh\). Небольшое число
параметров обеспечивает высокую скорость обучения, однако ограничивает
выразительную способность модели. Обучение проводилось в два этапа по 5000
эпох с шагами обучения \(10^{-4}\) и \(5\cdot10^{-5}\); суммарное время
составило 1 ч 36 мин.

Относительные \(L_2\)-ошибки на тестовом радиусе \(R=0.30\) приведены в
таблице~\ref{tab:light_errors}.

\begin{table}[H]
    \centering
    \caption{Относительные \(L_2\)-ошибки лёгкой архитектуры на тестовом радиусе \(R=0.30\).}
    \label{tab:light_errors}
    \begin{tabular}{c|cccc}
        \hline
        Поле & \(u\) & \(v\) & \(p\) & \(\alpha\) \\
        \hline
        \(L_2\), \% & 29.946 & 20.397 & 1.202 & 10.050 \\
        \hline
    \end{tabular}
\end{table}

Ошибки по скоростям достигают 20--30\%, по объемной доле --- около 10\%.
Давление восстанавливается точно благодаря преимущественно гидростатическому
характеру его распределения, которое слабо зависит от начального радиуса.
Полученные значения показывают, что четырёх слоёв по 128 нейронов недостаточно
для надёжного описания параметрической зависимости скоростного поля и
интерфейса.

\subsection{Средняя архитектура}

Средняя конфигурация содержит шесть скрытых слоёв по 256 нейронов и занимает
промежуточное положение по ёмкости между лёгкой и тяжёлой. Обучение проводилось
по той же двухстадийной схеме с теми же шагами обучения; время обучения
составило 3 ч 56 мин.

Результаты на тестовом радиусе \(R=0.30\) приведены в
таблице~\ref{tab:medium_errors}.

\begin{table}[H]
    \centering
    \caption{Относительные \(L_2\)-ошибки средней архитектуры на тестовом радиусе \(R=0.30\).}
    \label{tab:medium_errors}
    \begin{tabular}{c|cccc}
        \hline
        Поле & \(u\) & \(v\) & \(p\) & \(\alpha\) \\
        \hline
        \(L_2\), \% & 20.928 & 13.755 & 1.189 & 6.661 \\
        \hline
    \end{tabular}
\end{table}

По сравнению с лёгкой ошибка по \(u\) снижается с 29.95\% до 20.93\%, по
\(\alpha\) --- с 10.05\% до 6.66\%. Увеличение глубины и ширины сети даёт
заметный выигрыш по точности, однако ошибки по скоростям остаются
существенными, что указывает на необходимость ещё большей ёмкости модели
для данной задачи.

\subsection{Тяжёлая архитектура}

Тяжёлая архитектура --- наиболее глубокая из рассматриваемых: восемь скрытых
слоёв по 350 нейронов с функцией активации \(\tanh\). Время обучения в две
стадии составило 9 ч 40 мин.

В отличие от менее ёмких конфигураций, для тяжёлой архитектуры проверка
проводилась на двух тестовых значениях радиуса пузырька. Результаты приведены в
таблице~\ref{tab:hard_errors}.

\begin{table}[H]
    \centering
    \caption{Относительные \(L_2\)-ошибки (\%) тяжёлой архитектуры на двух тестовых значениях радиуса пузырька.}
    \label{tab:hard_errors}
    \begin{tabular}{l|cccc}
        \hline
        Тестовый радиус & \(u\) & \(v\) & \(p\) & \(\alpha\) \\
        \hline
        \(R=0.20\) & 6.287 & 5.594 & 0.667 & 7.142 \\
        \(R=0.30\) & 11.093 & 7.179 & 1.201 & 5.867 \\
        \hline
    \end{tabular}
\end{table}

Тяжёлая архитектура даёт наименьшие ошибки среди всех рассматриваемых
конфигураций. Ошибка по \(u\) при \(R=0.30\) составляет 11.09\% против
20.93\% у средней и 29.95\% у лёгкой. Для радиуса \(R=0.20\) ошибки
заметно ниже: 6.29\% по \(u\) и 7.14\% по \(\alpha\). Различие между
двумя тестовыми радиусами объясняется тем, что \(R=0.20\) окружён с
двух сторон обучающими значениями (\(0.18\) и \(0.22\)), тогда как
\(R=0.30\) расположен ближе к краю диапазона.

На рисунках~\ref{fig:param_r020} и~\ref{fig:param_r030_hard} показаны поля,
восстановленные тяжёлой архитектурой для двух тестовых значений радиуса пузырька в момент
времени \(t=3.0\).

\begin{figure}[H]
    \centering
    \includegraphics[width=\textwidth]{hard_2stage_R0.80_t3.000.png}
    \caption{Параметрическая PINN (тяжёлая), \(R=0.20\) (тестовый), \(t=3.0\).
    Слева~--- CFD, в центре~--- PINN, справа~--- разность. Ошибка
    сосредоточена вблизи интерфейса.}
    \label{fig:param_r020}
\end{figure}

\begin{figure}[H]
    \centering
    \includegraphics[width=\textwidth]{hard_2stage_R1.20_t3.000.png}
    \caption{Параметрическая PINN (тяжёлая), \(R=0.30\) (тестовый), \(t=3.0\).
    Модель корректно восстанавливает положение интерфейса и скоростное поле.}
    \label{fig:param_r030_hard}
\end{figure}

По совокупности метрик тяжёлая архитектура выбирается как основная для
последующих экспериментов.

\subsection{Итоговое сравнение архитектур}

В таблицах~\ref{tab:architecture_errors} и~\ref{tab:architecture_training_time}
сведены результаты всех трёх архитектур на тестовом радиусе \(R=0.30\) для
удобства сравнения.

\begin{table}[H]
    \centering
    \caption{Сводная таблица \(L_2\)-ошибок трёх архитектур параметрической PINN на тестовом радиусе \(R=0.30\).}
    \label{tab:architecture_errors}
    \begin{tabular}{c|ccc}
        \hline
        Поле & тяжёлая & средняя & лёгкая \\
        \hline
        \(u\) & 11.093\% & 20.928\% & 29.946\% \\
        \(v\) & 7.179\% & 13.755\% & 20.397\% \\
        \(p\) & 1.201\% & 1.189\% & 1.202\% \\
        \(\alpha\) & 5.867\% & 6.661\% & 10.050\% \\
        \hline
    \end{tabular}
\end{table}

\begin{table}[H]
    \centering
    \caption{Вычислительная стоимость трёх архитектур. Время вывода измерялось на 100 снимках;
             время CFD-расчёта тех же 100 снимков в Basilisk составляет 11.7 с.}
    \label{tab:architecture_training_time}
    \begin{tabular}{l|ccc}
        \hline
        & тяжёлая & средняя & лёгкая \\
        \hline
        Время обучения               & 9 ч 40 мин & 3 ч 56 мин & 1 ч 36 мин \\
        Вывод (100 снимков)          & 6.5 с      & 2.4 с      & 1.4 с      \\
        Ускорение относительно CFD   & $\times$1.8 & $\times$4.9 & $\times$8.4 \\
        \hline
    \end{tabular}
\end{table}

Качество монотонно возрастает с увеличением ёмкости сети. Переход от тяжёлой к
лёгкой увеличивает ошибку по \(u\) почти в 2.7 раза, по \(\alpha\) --- в 1.7
раза. Для давления разница между моделями несущественна.

По времени вывода картина обратная: лёгкая архитектура в 4.6 раза быстрее
тяжёлой (1.4 с против 6.5 с на 100 снимков) и в 8.4 раза быстрее
непосредственного CFD-расчёта. Таким образом, выбор архитектуры определяется
требуемым балансом между точностью и скоростью вывода: если приоритет ---
минимальная ошибка, предпочтительна тяжёлая; если важна максимальная скорость
вывода при допустимой погрешности --- лёгкая. Дальнейшие эксперименты
проводятся с тяжёлой архитектурой как наиболее точной.

\subsection{Влияние числа стадий обучения}

После выбора тяжёлой архитектуры проводится эксперимент с удлинённым расписанием
обучения. Базовая модель обучается в две стадии (10\,000 эпох суммарно,
время обучения 9 ч 40 мин); расширенная --- в четыре стадии с более плавным
снижением скорости обучения (20\,000 эпох суммарно, время обучения 19 ч 17 мин):
\[
\begin{array}{c|cccc}
    \hline
    \text{Стадия} & 1 & 2 & 3 & 4 \\
    \hline
    \text{Эпохи} & 5000 & 5000 & 5000 & 5000 \\
    \text{Шаг обучения} & 10^{-4} & 5\cdot10^{-5} & 10^{-5} & 5\cdot10^{-6} \\
    \hline
\end{array}
\]
Результаты сравнения для двух рассматриваемых расписаний обучения приведены в
таблице~\ref{tab:stages}.

\begin{table}[H]
    \centering
    \caption{Относительные \(L_2\)-ошибки (\%) при 2 и 4 стадиях обучения на двух тестовых значениях радиуса пузырька.}
    \label{tab:stages}
    \begin{tabular}{l|cccc|cccc}
        \hline
        & \multicolumn{4}{c|}{\(R=0.20\)} & \multicolumn{4}{c}{\(R=0.30\)} \\
        Модель & \(u\) & \(v\) & \(p\) & \(\alpha\) & \(u\) & \(v\) & \(p\) & \(\alpha\) \\
        \hline
        hard\_2stage & 6.287 & 5.594 & 0.667 & 7.142 & 11.093 & 7.179 & 1.201 & 5.867 \\
        hard\_4stage & 5.578 & 4.422 & 0.623 & 6.553 & 10.252 & 6.159 & 1.160 & 5.017 \\
        \hline
    \end{tabular}
\end{table}

Увеличение числа стадий с 2 до 4 даёт умеренное улучшение: ошибки по
скоростям снижаются примерно на 10--15\%, по объемной доле --- на 8--15\%.
Наиболее заметный эффект наблюдается для \(R=0.30\), где ошибка по \(\alpha\)
снизилась с 5.87\% до 5.02\%. Это свидетельствует о том, что тяжёлая архитектура
ещё не исчерпала свои возможности при 2 стадиях и дополнительные эпохи с
меньшим шагом обучения позволяют точнее согласовать данные с физическими
ограничениями.

На рисунке~\ref{fig:training_history} показаны кривые обучения всех
параметрических моделей. Более глубокие архитектуры сходятся медленнее,
однако достигают меньшего финального значения функции потерь. Переходы между
стадиями соответствуют моментам снижения скорости обучения и проявляются
как характерные перегибы на кривых.

\begin{figure}[H]
    \centering
    \includegraphics[width=\textwidth]{training_history_all.png}
    \caption{Кривые обучения параметрических PINN-моделей: зависимость
    суммарной функции потерь от номера эпохи для лёгкой, средней
    и тяжёлой архитектур. Вертикальные отметки соответствуют переходам между стадиями.}
    \label{fig:training_history}
\end{figure}

\subsection{Исследование влияния числа обучающих точек}

В базовой конфигурации для каждого временного слоя и каждого радиуса берётся:
\begin{itemize}
    \item \(\mathcal{D}_\alpha\): 500 точек у интерфейса + 400 точек в объёме,
    \item \(\mathcal{D}_f\): 400 точек у интерфейса (ближняя зона) + 2000 (дальняя зона интерфейса) + 3000 в объёме.
\end{itemize}

Проводится исследование, в котором поочерёдно и совместно удваивается количество
точек в \(\mathcal{D}_\alpha\) и \(\mathcal{D}_f\). Архитектура во всех
вариантах --- тяжёлая, расписание обучения --- 2 стадии. Результаты приведены в
таблице~\ref{tab:ablation}.

\begin{table}[H]
    \centering
    \caption{Исследование влияния числа обучающих точек: относительные \(L_2\)-ошибки (\%) на двух тестовых значениях радиуса пузырька.}
    \label{tab:ablation}
    \begin{tabular}{l|cccc|cccc}
        \hline
        & \multicolumn{4}{c|}{\(R=0.20\)} & \multicolumn{4}{c}{\(R=0.30\)} \\
        Конфигурация & \(u\) & \(v\) & \(p\) & \(\alpha\) & \(u\) & \(v\) & \(p\) & \(\alpha\) \\
        \hline
        базовая                                     & 6.287 & 5.594 & 0.667 & 7.142 & 11.093 & 7.179 & 1.201 & 5.867 \\
        \(\alpha\)$\times$2                         & 6.298 & 4.826 & 0.623 & 7.149 & 11.454 & 7.343 & 1.162 & 5.833 \\
        PDE$\times$2                                & 6.386 & 5.071 & 0.621 & 7.167 & 10.197 & 6.190 & 1.147 & 5.786 \\
        \(\alpha\)$\times$2 + PDE$\times$2          & 6.045 & 4.268 & 0.614 & 7.357 & 10.518 & 6.500 & 1.122 & 5.771 \\
        \hline
    \end{tabular}
\end{table}

Время обучения для каждой исследуемой конфигурации приведено в
таблице~\ref{tab:ablation_time}. Удвоение числа точек примерно вдвое увеличивает
время одной эпохи, поэтому все конфигурации с увеличенной выборкой обучаются
заметно дольше базовой модели.

\begin{table}[H]
    \centering
    \caption{Время обучения исследуемых конфигураций (тяжёлая архитектура, 2 стадии).}
    \label{tab:ablation_time}
    \begin{tabular}{l|c}
        \hline
        Конфигурация & Время обучения \\
        \hline
        базовая                             & 9 ч 40 мин \\
        \(\alpha\)$\times$2                 & 17 ч 24 мин \\
        PDE$\times$2                        & 17 ч 30 мин \\
        \(\alpha\)$\times$2 + PDE$\times$2  & 18 ч 40 мин \\
        \hline
    \end{tabular}
\end{table}

Увеличение числа PDE-точек (PDE$\times$2) даёт заметный выигрыш по скоростям
для \(R=0.30\): ошибка по \(u\) снижается с 11.09\% до 10.20\%, по \(v\) ---
с 7.18\% до 6.19\%. Увеличение числа \(\alpha\)-точек практически не влияет на
поля скорости, но незначительно улучшает \(v\) при \(R=0.20\). Совместное
удвоение двух типов точек не даёт стабильного выигрыша: ошибка по \(\alpha\)
при \(R=0.20\) при этом незначительно возрастает. Таким образом, наиболее
эффективным отдельным изменением является увеличение числа PDE-точек.

\subsection{Итоговое сравнение с простой PINN}

В таблице~\ref{tab:simple_vs_param} приведено сравнение лучшей параметрической
модели (hard\_2stage) с простой PINN на тестовом радиусе \(R=0.30\).

\begin{table}[H]
    \centering
    \caption{Сравнение простой и параметрической PINN на тестовом радиусе \(R=0.30\).}
    \label{tab:simple_vs_param}
    \begin{tabularx}{\textwidth}{X|cccc}
        \hline
        Модель & \(u\) & \(v\) & \(p\) & \(\alpha\) \\
        \hline
        Простая PINN (обучение \(R=0.25\), тест \(R=0.30\)) & 54.43\% & 29.53\% & 1.23\% & 54.96\% \\
        Параметрическая PINN hard\_2stage (\(R=0.30\))        & 11.09\% & 7.18\% & 1.20\% & 5.87\% \\
        \hline
    \end{tabularx}
\end{table}

Параметрическая PINN снижает ошибку по \(\alpha\) почти в 10 раз, по \(u\) ---
почти в 5 раз. Давление восстанавливается одинаково хорошо двумя моделями.
Разница по скоростям и объемной доле демонстрирует принципиальное преимущество
параметрической архитектуры: одна обученная модель корректно описывает семейство
режимов при изменении начального радиуса внутри изученного диапазона
(см.~рисунки~\ref{fig:param_r020} и~\ref{fig:param_r030_hard}).

Помимо точности, параметрическая PINN обладает практическим преимуществом по
скорости вывода. После однократного обучения (9 ч 40 мин для тяжёлой архитектуры) предсказание
100 пространственно-временных снимков занимает 6.5 с, тогда как получение тех
же 100 снимков путём прямого CFD-расчёта в Basilisk требует 11.7 с. Ключевое
отличие, однако, не в кратности ускорения (×1.8 для тяжёлой, ×8.4 для лёгкой), а в
природе вычисления: каждый новый радиус потребовал бы отдельного CFD-запуска,
тогда как обученная параметрическая сеть покрывает весь диапазон
$R \in [0.18, 0.32]$ без какого-либо переобучения. Это делает её пригодной в
качестве суррогатной модели для задач параметрического анализа и оптимизации.
